\documentclass[10pt,a4paper]{amsart}
\usepackage[margin=19mm]{geometry}
\usepackage{amsmath,amssymb,mathtools}
\usepackage[scr=rsfs,cal=euler]{mathalfa}
\usepackage{xfrac}
\usepackage[unicode,hidelinks,bookmarks=false]{hyperref}
\newcommand{\ie}{i.e.\ }
\newcommand{\cf}{cf.\ }
\newcommand{\resp}{respectively\ }
\newcommand{\Subsection}{\subsection}
\providecommand{\nobreakdash}{\nobreak\hbox{-}}
\setlength{\emergencystretch}{2em}
\multlinegap0pt
\usepackage{bm}
\usepackage{bbm}
\usepackage{dsfont}
\usepackage{xspace}
\usepackage{cancel}
\usepackage{mathtools}
\usepackage{amsthm}
\makeatletter
\@ifundefined{theorem}{\newtheorem{theorem}{Theorem}}{}
\@ifundefined{lemma}{\newtheorem{lemma}[theorem]{Lemma}}{}
\makeatother
\newcommand{\C}{\mathds{C}}
\title[Partition functions of two-dimensional Coulomb gases with ci]{Partition functions of two-dimensional Coulomb gases with circular root- and jump-type singularities}
\author{Kohei Noda}
\date{Source: arXiv:2510.00843v1 (CC BY 4.0)}
\begin{document}
\maketitle

\noindent\emph{Source and licence.} Partition functions of two-dimensional Coulomb gases with circular root- and jump-type singularities. Version arXiv:2510.00843v1 (CC BY 4.0), distributed under the Creative Commons Attribution 4.0 International licence, as recorded in the arXiv licence metadata for this exact version and shown on the abstract page https://arxiv.org/abs/2510.00843v1. Full rights notice: licenses/arxiv-2510.00843-CC-BY-4.0.txt.\par\smallskip

\noindent\textit{Editorial scope.} Counted unit: the Lemma whose source label is \texttt{lemma:S3} in the article (source0.tex lines 1302--1354, source label \texttt{lemma:S3}), delivered together with the article's own complete original proof of it (source0.tex lines 1356--1537, one \texttt{proof} environment of 182 lines). No mathematics is written, repaired or re-derived: statement and proof are the article's own text line for line. Required context retained uncounted: def of rtau ODE (lines 757--760); lemma:hnj asymptotics (lines 881--896); def of rtau g1plus (lines 920--937); lemma:S1 (lines 938--996); lemma:Euler-Maclaurin (lines 3022--3042). Every definition, display and label of those lines is retained unchanged. Omitted from this package and not redistributed: 1--756, 761--880, 897--919, 997--1301, 1355--1355, 1538--3021, 3043--3413. Nothing inside a retained excerpt is omitted or altered: every retained body is delivered exactly as the article's lines are written, so no omission receipt of any kind accompanies this package. Citations: the retained excerpts carry no citation command at all, so no bibliography entry of the article is transcribed and no reference is redistributed. Reference adaptation: none was needed and none was made; every reference inside the delivered unit resolves inside it, so no unresolved reference or citation is produced. Printed slips kept literal: no printed word, symbol, hypothesis or number of the counted statement or of its proof was altered. Layout and class: the article's own class is \texttt{amsart} with packages including geometry, amssymb, amsthm, amsmath, amsxtra, dsfont, appendix, stix, amsaddr, xcolor, hyperref, cleveref, todonotes, xparse; the wrapper is the lane's pinned \texttt{amsart} editorial wrapper at 10pt on A4, which restates the theorem-like environments the retained text needs and adds the article's own macro definitions that the retained text needs (\texttt{\textbackslash C}), with extra wrapper packages (bm, bbm, dsfont, xspace, cancel, mathtools). The delivered numbering is the unit's own. Assets: no retained span contains a figure environment, a graphics-inclusion command, a PDF-page inclusion, a table or a TikZ picture; no image file is copied into this package, the package declares no asset dependency and no image file is redistributed with it. Licence: version arXiv:2510.00843v1 of this article is distributed under the Creative Commons Attribution 4.0 International licence, as recorded in the arXiv licence metadata for this exact version and shown on the abstract page https://arxiv.org/abs/2510.00843v1. A full rights notice is delivered at licenses/arxiv-2510.00843-CC-BY-4.0.txt. Characters: every retained excerpt is pure ASCII, so the delivered engine, XeLaTeX with its default Latin Modern text font, reports no missing character.
\begin{equation}
\label{def of rtau ODE}
\frac{dr_{\tau}}{d\tau}=\frac{1}{2r_{\tau}\Delta Q(r_{\tau})}>0,
\end{equation}

\begin{lemma}\label{lemma:hnj asymptotics}
For $D_n\leq j\leq n-1$ and $\mathsf{k}(r)=2\alpha\log r$, as $n\to+\infty$, we have 
\begin{equation}
\label{def of Dn leq j leq n-1}
h_{n,j}
=\sqrt{\frac{2\pi}{n}} \frac{r_{\tau}}{\sqrt{\Delta Q(r_{\tau})}}e^{\mathsf{k}(r_{\tau})}e^{-nV_{\tau}(r_{\tau})}\cdot\Bigl(1+\frac{\mathcal{A}(r_{\tau})}{n}+\mathcal{O}(\frac{(\log n)^{\nu}}{j^{3/2}})\Bigr)    
\end{equation}
for some $\nu>0$, where
\begin{align*}
\mathcal{A}(r_{\tau})&:=\mathcal{B}(r_{\tau})+\frac{\mathsf{k}'(r_{\tau})^2}{2}\frac{1}{d_2}+\frac{\mathsf{k}''(r_{\tau})}{2}\frac{1}{d_2}+\frac{\mathsf{k}'(r_{\tau})}{r_{\tau}}\frac{1}{d_2}-\frac{\mathsf{k}'(r_{\tau})}{2}\frac{d_3}{d_2^2},   
\\
\mathcal{B}(r)&:=-\frac{1}{32}\frac{\partial_r^2\Delta Q(r)}{(\Delta Q(r))^2}-\frac{19}{96}\frac{\partial_r\Delta Q(r)}{(\Delta Q(r))^2}+\frac{5}{96}\frac{(\partial_r \Delta Q(r))^2}{(\Delta Q(r))^3}+\frac{1}{12}\frac{1}{r^2\Delta Q(r)},      \\
d_m&:=V_{\tau}^{(m)}(r_{\tau}).
\end{align*}
Here, $\mathcal{O}(j^{-3/2}(\log n)^{\nu})$ can be replaced with $\mathcal{O}(n^{-2})$ for large $j$, i.e., for $j\geq c_0 n$ with some $c_0>0$. 
\end{lemma}

\begin{align}
\label{def of rtau g1plus}
    r_{\tau(g_{1,+})}
&=
\rho+\frac{1}{2\sqrt{\Delta Q(\rho)}}\frac{\Delta_{n,+}}{\sqrt{n}}
-\frac{\Delta Q(\rho)+\rho\partial_r\Delta Q(\rho)}{8\rho \Delta Q(\rho)^2}\frac{\Delta_{n,+}^2}{n}+\mathcal{O}\Bigl(\frac{\Delta_{n,+}^3}{n^{3/2}}\Bigr), 
\\
\label{def of rtau g1minus}
r_{\tau(g_{1,-})}
&=\rho-\frac{1}{2\sqrt{\Delta Q(\rho)}}\frac{\Delta_{n,-}}{\sqrt{n}}
-\frac{\Delta Q(\rho)+\rho\partial_r\Delta Q(\rho)}{8\rho \Delta Q(\rho)^2}\frac{\Delta_{n,-}^2}{n}+\mathcal{O}\Bigl(\frac{\Delta_{n,-}^3}{n^{3/2}}\Bigr).
%\\
%\label{def of rtau j1minus}
%r_{\tau(j_{1,-})}
%&=
%\rho-\frac{1}{2\sqrt{\Delta Q(\rho)}}\frac{\Delta_{n,-}^{(\epsilon)}}{\sqrt{n}}
%-\frac{\Delta Q(\rho)+\rho\partial_r\Delta Q(\rho)}{8\rho \Delta Q(\rho)^2}\frac{{\Delta_{n,-}^{(\epsilon)}}^2}{n}+\mathcal{O}\Bigl(\frac{{\Delta_{n,-}^{(\epsilon)}}^3}{n^{3/2}}\Bigr),
\end{align}

\begin{lemma}\label{lemma:S1}
There exists $\delta>0$ such that as $n\to+\infty$, we have 
\begin{equation}
\label{def of S1 asymptotics}
S_1= C_1^{(1)}n + C_2^{(1)}\sqrt{n}+C_3^{(1)}+C_n^{(1)}+
\mathcal{O}\Bigl(\frac{(\log n)^{3}}{n^{\frac{1}{12}}}\Bigr),    
\end{equation}
uniformly for $u\in\{z\in\C:|z-x|\leq \delta\}$ and $a$ in compact subsets of $(-1,+\infty)$, where 
\begin{align*}
C_1^{(1)}&:=
\tau_{\rho}u
+
\int_{0}^{\rho}2a u\Delta Q(u)\log(\rho-u)\,du
\qquad
C_2^{(1)}:=0,
\\
C_3^{(1)}&:=
\frac{a}{2}\log 2
+\frac{a}{4}\log\Delta Q(\rho)
+\frac{a}{2}\log\rho
-
\frac{a(a-1)}{8}\Bigl(3+\frac{\rho\,\partial_r\Delta Q(\rho)}{\Delta Q(\rho)}\Bigr)
\\
&\quad
+
\frac{a}{4}\int_0^{\rho}\Bigl(\frac{1}{\rho-u} 
\frac{u\partial_r\Delta Q(u)}{\Delta Q(u)}
-
\frac{1}{\rho-u} 
\frac{\rho\partial_r\Delta Q(\rho)}{\Delta Q(\rho)}
\Bigr)\,du
\\
&\quad
-\frac{a}{4}
\log(2\rho\sqrt{\Delta Q(\rho)})
\Bigl(4\alpha+a+2-\frac{\rho\,\partial_r\Delta Q(\rho)}{\Delta Q(\rho)}\Bigr),
\\
C_n^{(1)}&:=
-\rho\sqrt{\Delta Q(\rho)}\sqrt{n}\Delta_{n,-}u
-D_nu
-aD_{n}\log\rho
\\
&\quad
-\frac{a}{2}\rho\sqrt{\Delta Q(\rho)}
\Bigl(2\log\Delta_{n,-}-2\log 2 -\log\Delta Q(\rho)-\log n-2\Bigr)\Delta_{n,-}\sqrt{n}
\\
&\quad
-\frac{a}{8}\Bigl(1+\frac{\rho\partial_r\Delta Q(\rho)}{\Delta Q(\rho)}\Bigr)\Delta_{n,-}^2
-\frac{a}{2}\log\Delta_{n,-}
+\frac{a}{4}\log n
+\frac{a(a-1)}{2}
\frac{\sqrt{n}}{\Delta_{n,-}}\rho\sqrt{\Delta Q(\rho)}
\\
&\quad
-\frac{a(a-1)(2a-3)}{12}\rho\sqrt{\Delta Q(\rho)}\frac{\sqrt{n}}{\Delta_{n,-}^3}
+\frac{a}{4}\Bigl(\log\Delta_{n,-}-\frac{1}{2}\log n\Bigr)
\Bigl( 4\alpha+a+2 -\frac{\rho\partial_r\Delta Q(\rho)}{\Delta Q(\rho)} \Bigr).
\end{align*}
\end{lemma}

\begin{lemma}[Euler-Maclaurin formula]\label{lemma:Euler-Maclaurin}
Let $f(x)$ be $2m$ times differentiable function on the interval $[p,q]$. Then we have 
\begin{align}
\begin{split}
\sum_{j=p+1}^{q-1}f(j)&=\int_{p}^{q}f(x)\,dx-\frac{f(p)+f(q)}{2}+\sum_{j=1}^{m-1}\frac{B_{2j}}{(2j)!}\Bigl( 
f^{(2j-1)}(q)-f^{(2j-1)}(p)
\Bigr)
+R_{2m}, 
\\
\sum_{j=p}^{q}f(j)&=\int_{p}^{q}f(x)\,dx+\frac{f(p)+f(q)}{2}+\sum_{j=1}^{m-1}\frac{B_{2j}}{(2j)!}\Bigl( 
f^{(2j-1)}(q)-f^{(2j-1)}(p)
\Bigr)
+R_{2m},
\\
\sum_{j=p}^{q-1}f(j)
&=\int_{p}^{q}f(x)\,dx-\frac{f(q)-f(p)}{2}+\sum_{j=1}^{m-1}\frac{B_{2j}}{(2j)!}\Bigl(f^{(2j-1)}(q)-f^{(2j-1)}(p)\Bigr)+R_{2m}, 
\end{split}    
\end{align}
where the sequence $\{B_{2k}\}_k$ are even indexed Bernoulli numbers $B_{2}=\frac{1}{6}$, $B_4=-\frac{1}{30},\dots$, and the remainder term $R_{2m}$ satisfies the bound $|R_{2m}|\leq c_{2m}\int_{p}^{q}|f^{(2m)}(x)|\,dx$ with $c_{2m}=\frac{2\zeta(2m)}{(2\pi)^{2m}}$. 
Here, $\zeta(s)$ is the Riemann zeta function. 
\end{lemma}

\begin{lemma}\label{lemma:S3}
There exists $\delta>0$ such that as $n\to+\infty$, we have 
\begin{equation}
\label{def of S3 asymptotics}
 S_3=
 C_1^{(3)}n+C_2^{(3)}\sqrt{n}+C_3^{(3)}+C_n^{(3)}+\mathcal{O}\Bigl(\frac{(\log n)^{\frac{1}{4}}}{n^{\frac{1}{4}}}\Bigr),
\end{equation}
uniformly for $u\in\{z\in\C:|z-x|\leq \delta\}$ and $a$ in compact subsets of $(-1,+\infty)$, where 
\begin{align*}
C_1^{(3)}&:=
a\int_{\rho}^{r_1}\log (t-\rho)\,2t\Delta Q(t)\,dt,
\\
C_2^{(3)}&:=0, 
\\
C_3^{(3)}&:=
-\frac{a}{2}\log(r_1-\rho)
+\frac{a}{2}\log2
+\frac{a}{4}\log(\Delta Q(\rho))
+\frac{a(a-1)}{4}\bigg[\frac{1}{2}\Bigl(1+\frac{\rho\,\partial_r\Delta Q(\rho)}{\Delta Q(\rho)}\Bigr)-\frac{\rho}{r_1-\rho}\bigg]
\\
&\quad
+\frac{a}{4}\Bigl(4\alpha+a+2-\frac{\rho\,\partial_r\Delta Q(\rho)}{\Delta Q(\rho)}\Bigr)\log(2(r_1-\rho)\sqrt{\Delta Q(\rho)})
\\
&\quad
-\frac{a}{4}\int_{\rho}^{r_1}\Bigl(\frac{1}{t-\rho}\frac{t\,\partial_t\Delta Q(t)}{\Delta Q(t)}-\frac{1}{t-\rho}\frac{\rho\,\partial_r\Delta Q(\rho)}{\Delta Q(\rho)}\Bigr)\,dt, 
\\
\begin{split}
C_n^{(3)}&:=
-\frac{a}{2}\rho\sqrt{\Delta Q(\rho)}
\Bigl( 
2\log\Delta_{n,+}-2\log 2-\log\Delta Q(\rho)-\log n-2
\Bigr)\sqrt{n}\Delta_{n,+}
\\
&\quad
+\frac{a}{8}\Bigl(1+\frac{\rho\partial_r\Delta Q(\rho)}{\Delta Q(\rho)}\Bigr)\Delta_{n,+}^2
-\frac{a}{2}\log(\Delta_{n,+})
+\frac{a}{4}\log n
+\frac{a(a-1)}{2}\rho\sqrt{\Delta Q(\rho)}\frac{\sqrt{n}}{\Delta_{n,+}}
\\
&\quad
-\frac{a}{4}
\Bigl( 
\log\Delta_{n,+}-\frac{1}{2}\log n
\Bigr)
\Bigl(
4\alpha+a+2-\frac{\rho\partial_r\Delta Q(\rho)}{\Delta Q(\rho)}
\Bigr)
-\frac{a(a-1)(2a-3)}{12}\rho\sqrt{\Delta Q(\rho)}\frac{\sqrt{n}}{\Delta_{n,+}^3}. 
\end{split}
\end{align*}


\end{lemma}

\begin{proof}
Similar manner to Lemma~\ref{lemma:S1}, there exists $c>0$ such that as $n\to+\infty$, we have  
\begin{align}
\begin{split}
\label{def of hnj out in S3}
h_{n,j}^{(\mathrm{out})}(\rho)
&=  
\frac{\sqrt{2\pi}r_{\tau}e^{\mathsf{k}(r_{\tau})}e^{-nV_{\tau}(r_{\tau})}}{\sqrt{n\Delta Q(r_{\tau})}}
(r_{\tau}-\rho)^a
\bigg[
1+\frac{1}{n}
\Bigl(
\mathcal{A}_1(r_{\tau})
+\mathcal{M}_1^{(\mathrm{out})}(r_{\tau})
\Bigr)
\\
&\quad
+\frac{1}{n^2}
\Bigl\{
\frac{a(a-1)(a-2)(a-3)u^4}{24d_2^2(r_{\tau}-\rho)^4}
+\mathcal{A}_2(r_{\tau})
+\mathcal{O}\Bigl(\frac{1}{(r_{\tau}-\rho)^3}\Bigr)
\Bigr\}
\bigg]
+
\mathcal{O}(e^{-c(\log n)^2}),
\end{split}
\end{align}
where the constant in the error term might depend on $a$, but the order of the error term is independent of $a$. Here, 
\[
\mathcal{M}_1^{(\mathrm{out})}(r_{\tau})
=
\frac{a(a-1)(r_{\tau}-\rho)^{-2}}{8\Delta Q(r_{\tau})}
+\frac{a(r_{\tau}-\rho)^{-1}}{8\Delta Q(r_{\tau})}
\Bigl(-\frac{\partial_r\Delta Q(r_{\tau})}{\Delta Q(r_{\tau})}
+\frac{4\alpha+3}{r_{\tau}}
\Bigr). 
\]
Also, there exists $c>0$ independent of $n$ such that we have 
\begin{equation}
\label{def of hnj in in S3}    
e^{u}\frac{h_{n,j}^{(\mathrm{in})}(\rho)}{h_{n,j}}=e^{-nV_{\tau}(r_{\tau})}\cdot\mathcal{O}(e^{-cn}), 
\end{equation}
where we can take the error term to be independent of $u$.
Thus, by Lemma~\ref{lemma:hnj asymptotics}, \eqref{def of hnj out in S3}, and \eqref{def of hnj in in S3}, let us choose $\delta>0$ sufficiently small so that 
\[
e^{u}\frac{h_{n,j}^{(\mathrm{in})}(\rho)}{h_{n,j}}+\frac{h_{n,j}^{(\mathrm{out})}(\rho)}{h_{n,j}}
\]
remains bounded away from the interval $(-\infty,0]$ as $n\to+\infty$ uniformly for $u\in\{z\in\C:|z-x|\leq \delta\}$ and $a$ in compact subsets of $(-1,+\infty)$. 
Therefore, as $n\to+\infty$, we have 
\begin{align*}
&\quad\sum_{j=g_{1,+}+1}^{n-1}\log\Bigl[
e^{u}\frac{h_{n,j}^{(\mathrm{in})}(\rho)}{h_{n,j}}+\frac{h_{n,j}^{(\mathrm{out})}(\rho)}{h_{n,j}}
 \Bigr]   
 \\
&= 
a\sum_{j=g_{1,+}+1}^{n-1}\log(r_{\tau}-\rho)
+
\frac{1}{n}
\sum_{j=g_{1,+}+1}^{n-1}
\Bigl\{ 
\frac{a(a-1)(r_{\tau}-\rho)^{-2}}{8\Delta Q(r_{\tau})}
+\frac{a(r_{\tau}-\rho)^{-1}}{8\Delta Q(r_{\tau})}
\Bigl(-\frac{\partial_r\Delta Q(r_{\tau})}{\Delta Q(r_{\tau})}
+\frac{4\alpha+3}{r_{\tau}}
\Bigr)
\Bigr\}
\\
&\quad
-\frac{1}{n^2}
\sum_{j=g_{1,+}+1}^{n-1}
\frac{a(a-1)(2a-3)}{64\Delta Q(r_{\tau})^2(r_{\tau}-\rho)^4}
+
\sum_{j=g_{1,+}+1}^{n-1}
\mathcal{O}\Bigl(\frac{(r_{\tau}-\rho)^{-3}}{n^2}\Bigr), 
\end{align*} 
uniformly for $u\in\{z\in\C:|z-x|\leq \delta\}$ and $a$ in compact subsets of $(-1,+\infty)$. 
By change of variable $r_{\tau(j)}=t$ and \eqref{def of rtau ODE}, we have 
\begin{equation}
\label{def of summation in 4 S3}
\sum_{j=g_{1,+}+1}^{n-1}
\mathcal{O}\Bigl(\frac{(r_{\tau}-\rho)^{-3}}{n^2}\Bigr)
=
\mathcal{O}\Bigl(\frac{(\log n)^{\frac{3}{8}}}{n^{\frac{7}{8}}}\Bigr). 
\end{equation}
By Lemma~\ref{lemma:Euler-Maclaurin}, we have 
\begin{align*}
\sum_{j=g_{1,+}+1}^{n-1}a\log(r_{\tau}-\rho)
&=  
\int_{g_{1,+}}^{n}a\log(r_{\tau(t)}-\rho)\,dt
-\frac{a}{2}\log(r_1-\rho)-\frac{a}{2}\log(r_{\tau(g_{1,+})}-\rho)
+\mathcal{O}(n^{-\frac{3}{8}}). 
\end{align*} 
Note that by \eqref{def of rtau g1plus}, as $n\to+\infty$, we have 
\begin{align*}
\int_{g_{1,+}}^{n}a\log(r_{\tau(t)}-\rho)\,dt
&=n\,a\int_{\rho}^{r_1}\log (u-\rho)\,2u\Delta Q(u)\,du
\\
&\quad
-a\rho\sqrt{\Delta Q(\rho)}
\Bigl( 
\log\Delta_{n,+}-\log 2-\frac{1}{2}\log\Delta Q(\rho)-\frac{1}{2}\log n-1
\Bigr)\sqrt{n}\Delta_{n,+}
\\
&\quad
+\frac{a}{8}\Bigl(1+\frac{\rho\partial_r\Delta Q(\rho)}{\Delta Q(\rho)}\Bigr)\Delta_{n,+}^2
+
\mathcal{O}\Bigl( 
\frac{\Delta_{n,+}^3}{\sqrt{n}}
\Bigr).
\end{align*}
Thus, the Taylor theorem together with \eqref{def of rtau g1plus} gives rise to
\begin{align}
\begin{split}
\label{def of summation in 1 S3}
\sum_{j=g_{1,+}+1}^{n-1}a\log(r_{\tau}-\rho)
&=  
n\,a\int_{r_{\tau}(g_{1,+})}^{r_1}\log (u-\rho)\,2u\Delta Q(u)\,du
\\
&\quad
+\frac{a}{2}\rho\sqrt{\Delta Q(\rho)}
\Bigl( 
2\log 2+2+\log\Delta Q(\rho)+\log n
-2\log\Delta_{n,+}
\Bigr)\sqrt{n}\Delta_{n,+}
\\
&\quad
+\frac{a}{8}\Bigl(1+\frac{\rho\partial_r\Delta Q(\rho)}{\Delta Q(\rho)}\Bigr)\Delta_{n,+}^2
-\frac{a}{2}\log(r_1-\rho)
-\frac{a}{2}\log(\Delta_{n,+})
+\frac{a}{2}\log2
\\
&\quad
+\frac{a}{4}\log(\Delta Q(\rho))
+\frac{a}{4}\log n
+
\mathcal{O}\Bigl( 
\frac{1}{\Delta_{n,+}}
\Bigr). 
\end{split}
\end{align}
The term of order $n^{-1}$ can be calculated as
\begin{align}
\begin{split}
\label{def of summation in 2 S3}
&\quad
\frac{1}{n}
\sum_{j=g_{1,+}+1}^{n-1}
\Bigl\{ 
\frac{a(a-1)(r_{\tau}-\rho)^{-2}}{8\Delta Q(r_{\tau})}
+\frac{a(r_{\tau}-\rho)^{-1}}{8\Delta Q(r_{\tau})}
\Bigl(-\frac{\partial_r\Delta Q(r_{\tau})}{\Delta Q(r_{\tau})}
+\frac{4\alpha+3}{r_{\tau}}
\Bigr)
\Bigr\}
\\
&=
\frac{a(a-1)}{4}\bigg[\frac{1}{2}\Bigl(1+\frac{\rho\,\partial_r\Delta Q(\rho)}{\Delta Q(\rho)}\Bigr)+\log(2(r_1-\rho)\sqrt{\Delta Q(\rho)})-\frac{\rho}{r_1-\rho}\bigg]
\\
&\quad
+\frac{a}{4}\Bigl(4\alpha+3-\frac{\rho\,\partial_r\Delta Q(\rho)}{\Delta Q(\rho)}\Bigr)\log(2(r_1-\rho)\sqrt{\Delta Q(\rho)})
\\
&\quad
-\frac{a}{4}\int_{\rho}^{r_1}\Bigl(\frac{1}{u-\rho}\frac{u\,\partial_u\Delta Q(u)}{\Delta Q(u)}-\frac{1}{u-\rho}\frac{\rho\,\partial_r\Delta Q(\rho)}{\Delta Q(\rho)}\Bigr)\,du
\\
&\quad
+\frac{a(a-1)}{2}\rho\sqrt{\Delta Q(\rho)}\frac{\sqrt{n}}{\Delta_{n,+}}
+\frac{a}{4}\Bigl(\frac{1}{2}\log n-\log\Delta_{n,+}\Bigr)\Bigl(4\alpha+a+2-\frac{\rho\,\partial_r\Delta Q(\rho)}{\Delta Q(\rho)}\Bigr)+\mathcal{O}(\Delta_{n,+}^{-2}). 
\end{split}
\end{align}
Finally, it is straightforward to see that by \eqref{def of rtau g1plus} and Lemma~\ref{lemma:Euler-Maclaurin}, 
\begin{align}
\label{def of summation in 3 S3}
-\frac{1}{n^2}
\sum_{j=g_{1,+}+1}^{n-1}
\frac{a(a-1)(2a-3)}{64\Delta Q(r_{\tau})^2(r_{\tau}-\rho)^4}
&=
-\frac{a(a-1)(2a-3)}{12}\rho\sqrt{\Delta Q(\rho)}\frac{\sqrt{n}}{\Delta_{n,+}^3}
+\mathcal{O}\Bigl(\frac{(\log n)^{1/4}}{n^{1/4}}\Bigr). 
\end{align}
Combining \eqref{def of summation in 4 S3} with \eqref{def of summation in 1 S3}, \eqref{def of summation in 2 S3}, and \eqref{def of summation in 3 S3}, we obtain \eqref{def of S3 asymptotics}.
\end{proof}

\end{document}
